\documentclass[10pt,a4paper]{amsart}
\usepackage{amsmath,amssymb,amsthm,graphicx}
% Retained minimal preamble for the Stacks Project chapter more-algebra.tex (More on Algebra).
% The chapter file pinned-original-source.tex carries only \input{preamble} on its first line, so no
% preamble text lives inside the chapter itself.  The Stacks Project uses one single shared preamble for
% every chapter; the snapshot stages it identically as sources/stacks-chow/preamble.tex and
% sources/stacks-etale-cohomology/preamble.tex.  The block below is copied mechanically, line for line,
% from sources/stacks-chow/preamble.tex lines 14-18 (the Xy-pic setup) and 154-217 (the native theorem
% declarations together with the mathematical macros).  Nothing was added, removed or redefined; the
% wrapper class and the amsmath/amssymb/amsthm/graphicx base packages are supplied by the assembler.
% Chapter: more-algebra.tex -- More on Algebra.
% For commutative diagrams we use Xy-pic
\usepackage[all]{xy}

% We use 2cell for 2-commutative diagrams.
\xyoption{2cell}

% Theorem environments.
%
\theoremstyle{plain}
\newtheorem{theorem}[subsection]{Theorem}
\newtheorem{proposition}[subsection]{Proposition}
\newtheorem{lemma}[subsection]{Lemma}

\theoremstyle{definition}
\newtheorem{definition}[subsection]{Definition}
\newtheorem{example}[subsection]{Example}
\newtheorem{exercise}[subsection]{Exercise}
\newtheorem{situation}[subsection]{Situation}

\theoremstyle{remark}
\newtheorem{remark}[subsection]{Remark}
\newtheorem{remarks}[subsection]{Remarks}

\numberwithin{equation}{subsection}

% Macros
%
\def\lim{\mathop{\mathrm{lim}}\nolimits}
\def\colim{\mathop{\mathrm{colim}}\nolimits}
\def\Spec{\mathop{\mathrm{Spec}}}
\def\Hom{\mathop{\mathrm{Hom}}\nolimits}
\def\Ext{\mathop{\mathrm{Ext}}\nolimits}
\def\SheafHom{\mathop{\mathcal{H}\!\mathit{om}}\nolimits}
\def\SheafExt{\mathop{\mathcal{E}\!\mathit{xt}}\nolimits}
\def\Sch{\mathit{Sch}}
\def\Mor{\mathop{\mathrm{Mor}}\nolimits}
\def\Ob{\mathop{\mathrm{Ob}}\nolimits}
\def\Sh{\mathop{\mathit{Sh}}\nolimits}
\def\NL{\mathop{N\!L}\nolimits}
\def\CH{\mathop{\mathrm{CH}}\nolimits}
\def\proetale{{pro\text{-}\acute{e}tale}}
\def\etale{{\acute{e}tale}}
\def\QCoh{\mathit{QCoh}}
\def\Ker{\mathop{\mathrm{Ker}}}
\def\Im{\mathop{\mathrm{Im}}}
\def\Coker{\mathop{\mathrm{Coker}}}
\def\Coim{\mathop{\mathrm{Coim}}}

% Boxtimes
%
\DeclareMathSymbol{\boxtimes}{\mathbin}{AMSa}{"02}

%
% Macros for moduli stacks/spaces
%
\def\QCohstack{\mathcal{QC}\!\mathit{oh}}
\def\Cohstack{\mathcal{C}\!\mathit{oh}}
\def\Spacesstack{\mathcal{S}\!\mathit{paces}}
\def\Quotfunctor{\mathrm{Quot}}
\def\Hilbfunctor{\mathrm{Hilb}}
\def\Curvesstack{\mathcal{C}\!\mathit{urves}}
\def\Polarizedstack{\mathcal{P}\!\mathit{olarized}}
\def\Complexesstack{\mathcal{C}\!\mathit{omplexes}}
% \Pic is the operator that assigns to X its picard group, usage \Pic(X)
% \Picardstack_{X/B} denotes the Picard stack of X over B
% \Picardfunctor_{X/B} denotes the Picard functor of X over B
\def\Pic{\mathop{\mathrm{Pic}}\nolimits}
\def\Picardstack{\mathcal{P}\!\mathit{ic}}
\def\Picardfunctor{\mathrm{Pic}}
\def\Deformationcategory{\mathcal{D}\!\mathit{ef}}
\title{The inequality [Gamma\_B : Gamma\_A] times [kappa\_B : kappa\_A] at most [L : K] for a finite extension of valuation rings}
\author{The Stacks Project Authors}
\date{stacks-project:more-algebra.tex:local-snapshot-20261012; GNU FDL 1.2 or later}
\begin{document}
\maketitle
\setcounter{subsection}{1}
\noindent\textit{Editorial note.} All delivered mathematical text is the author's own, copied from the pinned Stacks Project chapter file more-algebra.tex; no statement or proof step was rewritten, reordered or removed. The counted statement is the labeled lemma environment at source lines 37333-37342 and the counted proof is the proof environment at source lines 37344-37373, which immediately follows it and is its own proof.\par\medskip
\noindent\textit{Editorial note.} The remark that precedes the lemma at source lines 37330-37331, recording that the extension of fraction fields is not required to be finite, is outside the counted spans because it belongs to the surrounding discussion and not to the lemma or its proof.\par\medskip
\noindent\textit{Editorial note.} The chapter file pinned-original-source.tex carries only an input of the shared preamble on its first line, so no preamble text lives inside the chapter. The retained file sources/stacks-more-algebra/retained-preamble.tex is a mechanical, line-for-line copy of the project's single shared preamble as staged in the snapshot (Xy-pic setup plus the native theorem declarations and mathematical macros). Nothing was added, removed or redefined, and the excerpt serves only as preamble, never as a span coordinate.\par\medskip
\noindent\textit{Editorial note.} The source URL is the source-level Stacks Project URL, as required when more than one card is produced from one source; the per-claim tag is carried by claim\_key. The counted statement and proof refer to no labelled item of the chapter and cite no bibliography entry, so this card contributes no external reference stub and needs no bibliography.\par\medskip
\noindent\textit{Allowed standard prerequisites (editorial).}\begin{itemize}
\item A finite family of elements of a valued field with pairwise distinct valuations cannot sum to zero: after a minimal valuation is normalized to zero, the remaining terms have strictly larger valuation while a unit term cannot be cancelled by them. used at proof lines 37353-37371, where the index set P of minimal valuation is introduced, the coefficients are normalized and the sum over P is shown to be nonzero. Supplies the contradiction that proves the products b\_i c\_j are K-linearly independent.
\item If b\_1, ..., b\_n are units whose residue classes in kappa\_B are linearly independent over kappa\_A and the coefficients a\_ij are units, then the sum of the a\_ij b\_i over the index set P has valuation zero, hence is a unit of B. used at proof lines 37362-37366. Identifies the valuation of the partial sum over P with the valuation of the chosen minimal term, which is what makes the whole sum nonzero.
\item Nonzero elements c\_1, ..., c\_m of B whose images in Gamma\_B over Gamma\_A are pairwise distinct force, for a fixed minimal valuation, the equality j = j\_0 of the second index. used at proof lines 37346-37361. Reduces the index set P of minimal valuation to a single value of j, which is what allows the residue classes of the b\_i to be used in the next step.
\end{itemize}\medskip
\begin{lemma}
\label{lemma-inequality-general}
Let $A \subset B$ be an extension of valuation rings with
fraction fields $K \subset L$. If the extension $L/K$
is finite, then the residue field extension is finite,
the index of $\Gamma_A$ in $\Gamma_B$ is finite, and
$$
[\Gamma_B : \Gamma_A] [\kappa_B : \kappa_A] \leq [L : K].
$$
\end{lemma}

\begin{proof}
Let $b_1, \ldots, b_n \in B$ be units whose images in $\kappa_B$
are linearly independent over $\kappa_A$. Let $c_1, \ldots, c_m \in B$
be nonzero elements whose images in $\Gamma_B/\Gamma_A$ are pairwise
distinct. We claim that $b_i c_j$ are $K$-linearly independent
in $L$. Namely, we claim a sum
$$
\sum a_{ij} b_i c_j
$$
with $a_{ij} \in K$ not all zero cannot be zero. Choose $(i_0, j_0)$ with
$v(a_{i_0j_0}b_{i_0}c_{j_0})$ minimal. Replace $a_{ij}$ by
$a_{ij}/a_{i_0j_0}$, so that $a_{i_0 j_0} = 1$. Let
$$
P = \{(i, j) \mid
v(a_{ij}b_ic_j) = v(a_{i_0j_0}b_{i_0}c_{j_0}) \}
$$
By our choice of $c_1, \ldots, c_m$ we see that $(i, j) \in P$ implies
$j = j_0$. Hence if $(i, j) \in P$, then $v(a_{ij}) = v(a_{i_0j_0}) = 0$,
i.e., $a_{ij}$ is a unit. By our choice of $b_1, \ldots, b_n$
we see that
$$
\sum\nolimits_{(i, j) \in P} a_{ij}b_i
$$
is a unit in $B$. Thus the valuation of
$\sum\nolimits_{(i, j) \in P} a_{ij}b_ic_j$ is
$v(c_{j_0}) = v(a_{i_0j_0}b_{i_0}c_{j_0})$.
Since the terms with $(i, j) \not \in P$ in the first displayed sum
have strictly bigger valuation, we conclude that this sum cannot be
zero, thereby proving the lemma.
\end{proof}

\end{document}
